\subsection{Задания для лабораторной работы}
Используйте асинхронные таймеры, глобальные переменные для таймеров и состояний, полный проход по индексам \verb|0..24|. Формулировки ниже предельно простые; рядом указаны подсказки.
\subsection{Базовые алгоритмы}

\begin{taskbox}[Одновременное включение всей светодиодной линейки зелёным цветом]
\indent Требуется реализовать алгоритм, обеспечивающий однократную установку зелёного цвета с компонентами \verb|{0,1,0}| на всех светодиодах линейки с индексами от \verb|0| до \verb|24|. Для этого необходимо создать объект светодиодной линейки, реализовать функцию массовой установки цвета и выполнить её вызов.

\noindent\textbf{Ожидаемый результат:} вся линейка равномерно светится зелёным цветом.
\end{taskbox}

\begin{taskbox}[Отложенное включение синего цвета]
\indent Необходимо реализовать алгоритм, при котором светодиодная линейка остаётся выключенной в течение первых трёх секунд работы программы, после чего все светодиоды одновременно загораются синим цветом \verb|{0,0,1}|. Реализация должна использовать механизм отложенного вызова без применения блокирующих задержек.

\noindent\textbf{Ожидаемый результат:} через 3 секунды после запуска программы линейка полностью загорается синим цветом.
\end{taskbox}

\begin{taskbox}[Периодическое мигание красным цветом]
\indent Требуется реализовать периодическое переключение всей светодиодной линейки между состояниями «включено красным цветом \verb|{1,0,0}|» и «выключено \verb|{0,0,0}|» с периодом 0{,}5~с. Алгоритм должен быть основан на использовании асинхронного таймера и логической переменной состояния.

\noindent\textbf{Ожидаемый результат:} равномерное мигание линейки с периодом 0{,}5~с.
\end{taskbox}

\begin{taskbox}[Эффект «бегущего огня» (один активный светодиод)]
\indent Необходимо реализовать алгоритм, при котором в каждый момент времени включён только один светодиод зелёного цвета \verb|{0,1,0}|, а остальные выключены. Активный светодиод должен последовательно перемещаться от индекса \verb|0| к \verb|24| с циклическим повторением.

\noindent\textbf{Ожидаемый результат:} визуальный эффект перемещающегося «огонька» по всей линейке.
\end{taskbox}

\begin{taskbox}[Модель светофора на основе конечного автомата]
\indent Требуется реализовать конечный автомат с тремя состояниями: красный \verb|{1,0,0}|, жёлтый \verb|{1,1,0}| и зелёный \verb|{0,1,0}|. Состояния должны циклически сменять друг друга с заданным временным интервалом.

\noindent\textbf{Ожидаемый результат:} периодическая и стабильная смена цветов по схеме «красный -- жёлтый -- зелёный».
\end{taskbox}
